\documentclass[a4paper, 10pt]{article}

% -- Coments
\usepackage{verbatim}

% ----- Fonts -----
% -- Fuente --
% \usepackage{fontspec}
% \setmonofont{JetBrainsMono Nerd Font}  

% -- Color --
\usepackage{xcolor}
%\definecolor{azul}{RGB}{00,33,99}
\definecolor{azul}{RGB}{35,72,180}

% -- Page Margin --
\usepackage[margin=1in]{geometry}

% -- Espaciados --
\newcommand{\Pspace}{0.5cm}
\newcommand{\Aspace}{0.2cm}

% -- Columnas --
\usepackage{multicol}

% -- Imagenes --
\usepackage{graphicx}
\usepackage{float}

% -- Matemáticas --
\usepackage{amsmath, amssymb}

% -- Gráficas --
\usepackage{pgfplots}
\pgfplotsset{compat=1.18}

% -- Código --
\usepackage{listings}
\lstset{
    language=Lisp,                   % Lenguaje del código
    basicstyle=\ttfamily\small,     % Fuente del código
    keywordstyle=\color{blue},      % Color de palabras clave
    commentstyle=\color{gray},      % Color de comentarios
    stringstyle=\color{red},        % Color de cadenas
    numbers=left,                   % Números de línea a la izquierda
    numberstyle=\tiny\color{gray},
    breaklines=true,                % Permitir saltos de línea
    frame=single                    % Marco alrededor del código
}

\begin{comment}
-- Formato de pregunta simple
    % - Problema n
    \vspace{\Pspace}
    \item Problema \par
        % Respuesta:
        \vspace{\Aspace} \par
        { \color{azul}  }


 -- Formato de pregunta multimple
        % - Problema 
        \vspace{\Pspace}
        \item  \par
             % Respuestas:
            \vspace{\Aspace} \par
            a) 
            \\ { \color{azul} }

            \vspace{\Aspace} \par
            b) 
            \\ { \color{azul} }


 -- Formato de imagen
\begin{figure}[!ht]
    \centering
    \includegraphics[width=0.5\textwidth]{}
\end{figure}


 -- Formato de tabla
\begin{tabular}{c|c}
    \textbf{A} & \textbf{B} \\
    \hline
    Texto & Texto \\
    Texto & Texto
\end{tabular} \\
\end{comment}


\title
{
    Lenguajes de Programación 2026-2 \\
    Examen Intermedio
}

\begin{document}

\maketitle

\begin{center}
    \begin{tabular}{r|l}
        \textbf{Expediente} & \textbf{Nombre} \\ \hline
        219208106 & Bórquez Guerrero Angel Fernando \\
    \end{tabular}
\end{center}

\rule{\linewidth}{0.3mm}

\vspace{0.3cm}
{\Large Grámatica $G$} \\
\begin{tabular}{rcll}
    Expression & ::= & Number & (const-exp $n$) \\
    & | & Identifier & (var-exp $x$) \\
    & | & zero?(Expression) & (zero?-exp $e_1$) \\
    & | & -(Expression, Expression) & (diff-exp $e_1$ $e_2$) \\
    & | & if Expression then Expression else Expression & (if-exp $e_1$ $e_2$ $e_3$) \\
    & | & let Identifier = Expression in Expression & (let-exp $x$ $e_1$ $e_2$) \\
    & | & proc(Identifier) Expression & (proc-exp $x$ $e_1$) \\
    & | & (Expression Expression) & (call-exp $e_1$ $e_2$) \\
\end{tabular} \\

\begin{enumerate}
    \item Escribe un programa del lenguaje PROC usando únicamente los constructos de $G$ tal que:
    \begin{itemize}
        \item Todos los constructos de $G$ aparecen en tu programa.
        \item El valor del programa es distinto con alcance léxico y con alcance dinámico.
    \end{itemize}
\begin{lstlisting}
let x = 60 in
    let fun = proc (y) -(y, x) in
        let x = 50 in
            if zero?(fun 60) then 1 else 0
\end{lstlisting}


    \item Escribe y justifica el valor del programa del problema 1.
    \begin{itemize}
        \item Con alcance léxico. \\
            {\color{azul} Se declara una variable $x$ con valor de 60, seguido de una función "fun(y) = y-x" que con alcance léxico se reemplaza esta $x$ por 60 quedando "fun(y) = y - 60", luego redeclaramos $x = 50$, creamos una expresión if que evalúa si fun(60) == 0 y dado que es cierto se devuelve 1.}
        \item Con alcance dinámico. \\
            {\color{azul} Al momento de llamar la función fun(60) se declaró $x = 50$, por lo tanto "fun(y) = y - 50", entonces nuestra expresión fun(60) == 0 será falsa retornando así un 0.}
    \end{itemize}



    \item Escribe el \textsl{AST} del programa del problema 1.
\begin{lstlisting}
(let-exp 'x (const-exp 60)
    (let-exp 'fun (proc-exp 'y (diff-exp (var-exp 'y) (var-exp 'x)))
        (let-exp 'x (const-exp 50)
            (if-exp (zero?-exp (call-exp (var-exp 'fun) (const-exp 60))
                    (const-exp 1)
                    (const-exp 0))))))
\end{lstlisting}
\end{enumerate}
\end{document}
